\begin{textsample}{NEG Dim 5 – global\_north\_2024\_10 – Score 1.00 – global\_north\_2024\_10/116882031.txt}  \label{ex:f5_neg_047}
Radiohead frontman Thom Yorke walked off stage during a solo show in Melbourne on Wednesday night after being heckled by a pro-Palestine protester in the crowd.
Footage from concert-goers captured a man in the crowd yelling at Yorke.
While it was difficult to hear his full comments, he said"the Israeli genocide of Gaza"and then referred to the death toll, saying that"half of them were children".
Yorke responded :"Come up here and say that.
Right here, come on.
Hop up on the fucking stage and say what you wan na say.
Don't stand there like a coward, come here and say it.
You want to piss on everybody ' s night?"The protester then yelled out :"How many dead children will it take for you to condemn the genocide in Gaza?"Yorke responded,"OK, you do it, see you later then", and walked offstage.
He returned a few minutes later to perform his final song of the evening, Radiohead ' s @ @ @ @ @ @ @ @ @ @ the end of the concert at the Sidney Myer Music Bowl, the second of two in Melbourne as part of Yorke ' s Everything tour, featuring music from across his career, including solo material and songs from Radiohead and the Smile.
He is scheduled to play the Sydney Opera House forecourt on Friday 1 November and Saturday 2 November.
Radiohead played Tel Aviv in 2017, defying a BDS-led call to boycott the country that included public criticism from figures including the British director Ken Loach.
In a statement on X at that time, responding directly to Loach, Yorke said :"Playing in a country isn't the same as endorsing the government.
We ' ve played in Israel for over 20 years through a succession of governments, some more liberal than others.
As we have in America.
We don't endorse Netanyahu any more than Trump, but we still play in America."Radiohead has a long history with Israel, with their early hit Creep first finding success on Israeli radio, and @ @ @ @ @ @ @ @ @ @.
But pressure on the band and its members to boycott Israel has grown over the past year.
In May, Radiohead and the Smile musician Jonny Greenwood was criticised for playing a gig in Tel Aviv with Israeli artist Dudu Tassa, with the BDS movement accusing him of"artwashing genocide".
Responding in a statement on his social media accounts, Greenwood, who is married to Israeli visual artist Sharona Katan and has collaborated with Israeli musicians previously, lamented"the silencing of this -- or any -- artistic effort made by Israeli Jews"."No art is as ' important ' as stopping all the death and suffering around us,"he said."How can it be?
But doing nothing seems like a worse option.
And silencing Israeli artists for being born Jewish in Israel doesn't seem like any way to reach an understanding between the two sides of this apparently endless conflict."

% matched lemmas: 
\end{textsample}
